\documentclass[10pt,letterpaper]{article}
%Mise en page
\usepackage[left=2cm, right=2cm, top=0.8in, bottom=0.7in]{geometry}
\usepackage{fancyhdr}
\usepackage{fancybox}
\usepackage{pdfpages} 
\renewcommand{\headrulewidth}{1.5pt}
\renewcommand{\footrulewidth}{1.5pt}
\pagestyle{fancy}
\newcommand\Loadedframemethod{TikZ}
\usepackage[framemethod=\Loadedframemethod]{mdframed}
\usepackage{tikz}
\usepackage[linesnumbered,ruled,vlined]{algorithm2e}
%\usepackage{url}
\usepackage{dsfont}
\usepackage{amssymb,amsmath}
\usepackage{xspace}

\lhead{
\textbf{New York University}
}
\rhead{\textbf{Fall 2026}
}
\chead{\textbf{
CSCI-GA 3033
 }}

\newcommand{\RR}{\mathds{R}}
\newcommand{\sign}{\mathop{\mathrm{sign}}}
\newcommand{\argmin}{\mathop{\mathrm{argmin}}}
\newcommand{\zero}{\mathbf{0}}
\newcommand{\one}{\mathbf{1}}
\newcommand{\bv}{\mathbf{b}}
\newcommand{\wv}{\mathbf{w}}
\newcommand{\xv}{\mathbf{x}}
\newcommand{\yv}{\mathbf{y}}
\newcommand{\zv}{\mathbf{z}}
\newcommand{\rv}{\mathbf{r}}
\newcommand{\inner}[2]{\langle #1, #2 \rangle}
\newcommand{\red}[1]{{\color{red}#1}}
\newcommand{\blue}[1]{{\color{blue}#1}}
\newcommand{\magenta}[1]{{\color{magenta}#1}}


\newcommand{\ea}{{et al.}\xspace}
\newcommand{\eg}{{e.g.}\xspace}
\newcommand{\ie}{{i.e.}\xspace}
\newcommand{\iid}{{i.i.d.}\xspace}
\newcommand{\cf}{{cf.}\xspace}
\newcommand{\wrt}{{w.r.t.}\xspace}
\newcommand{\aka}{{a.k.a.}\xspace}
\newcommand{\etc}{{etc.}\xspace}
\newcommand{\Dc}{\mathcal{D}}

\newcommand{\ans}[1]{{\color{orange}\textsf{Ans}: #1}}


\lfoot{}
\cfoot{\textbf{Gautam Kamath \textcopyright 2026}}
\rfoot{\thepage}

%================================
%================================

\setlength{\parskip}{1cm}
\setlength{\parindent}{1cm}
\tikzstyle{titregris} =
[draw=gray,fill=white, shading = exersicetitle, %
text=gray, rectangle, rounded corners, right,minimum height=.3cm]
\pgfdeclarehorizontalshading{exersicebackground}{100bp}
{color(0bp)=(green!40); color(100bp)=(black!5)}
\pgfdeclarehorizontalshading{exersicetitle}{100bp}
{color(0bp)=(red!40);color(100bp)=(black!5)}

\makeatletter
\def\mdf@@exercisepoints{}%new mdframed key:
\define@key{mdf}{exercisepoints}{%
\def\mdf@@exercisepoints{#1}
}

\mdfdefinestyle{exercisestyle}{%
outerlinewidth=1em,outerlinecolor=white,%
leftmargin=-1em,rightmargin=-1em,%
middlelinewidth=0.5pt,roundcorner=3pt,linecolor=black,
apptotikzsetting={\tikzset{mdfbackground/.append style ={%
shading = exersicebackground}}},
innertopmargin=0.1\baselineskip,
skipabove={\dimexpr0.1\baselineskip+0\topskip\relax},
skipbelow={-0.1em},
needspace=0.5\baselineskip,
frametitlefont=\sffamily\bfseries,
settings={\global\stepcounter{exercise}},
singleextra={%
\node[titregris,xshift=0.5cm] at (P-|O) %
{~\mdf@frametitlefont{\theexercise}~};
\ifdefempty{\mdf@@exercisepoints}%
{}%
{\node[titregris,left,xshift=-1cm] at (P)%
{~\mdf@frametitlefont{\mdf@@exercisepoints points}~};}%
},
firstextra={%
\node[titregris,xshift=1cm] at (P-|O) %
{~\mdf@frametitlefont{\theexercise}~};
\ifdefempty{\mdf@@exercisepoints}%
{}%
{\node[titregris,left,xshift=-1cm] at (P)%
{~\mdf@frametitlefont{\mdf@@exercisepoints points}~};}%
},
}
\makeatother


%%%%%%%%%

%%%%%%%%%%%%%%%
\mdfdefinestyle{theoremstyle}{%
outerlinewidth=0.01em,linecolor=black,middlelinewidth=0.5pt,%
frametitlerule=true,roundcorner=2pt,%
apptotikzsetting={\tikzset{mfframetitlebackground/.append style={%
shade,left color=white, right color=blue!20}}},
frametitlerulecolor=black,innertopmargin=1\baselineskip,%green!60,
innerbottommargin=0.5\baselineskip,
frametitlerulewidth=0.1pt,
innertopmargin=0.7\topskip,skipabove={\dimexpr0.2\baselineskip+0.1\topskip\relax},
frametitleaboveskip=1pt,
frametitlebelowskip=1pt
}
\setlength{\parskip}{0mm}
\setlength{\parindent}{10mm}
\mdtheorem[style=theoremstyle]{exercise}{\textbf{Exercise}}

%================Liste definition--numList-and alphList=============
\newcounter{alphListCounter}
\newenvironment
{alphList}
{\begin{list}
{\alph{alphListCounter})}
{\usecounter{alphListCounter}
\setlength{\rightmargin}{0cm}
\setlength{\leftmargin}{0.5cm}
\setlength{\itemsep}{0.2cm}
\setlength{\partopsep}{0cm}
\setlength{\parsep}{0cm}}
}
{\end{list}}
\newcounter{numListCounter}
\newenvironment
{numList}
{\begin{list}
{\arabic{numListCounter})}
{\usecounter{numListCounter}
\setlength{\rightmargin}{0cm}
\setlength{\leftmargin}{0.5cm}
\setlength{\itemsep}{0cm}
\setlength{\partopsep}{0cm}
\setlength{\parsep}{0cm}}
}
{\end{list}}

\usepackage[breaklinks=true,colorlinks=true,linkcolor=magenta,urlcolor=magenta,citecolor=black]{hyperref}

\usepackage{cleveref}
\usepackage{xpatch}
\xpretocmd{\algorithm}{\hsize=\linewidth}{}{}


%===========================================================

\usepackage{stmaryrd}
\usepackage{enumitem}
\newcommand{\pred}[1]{\llbracket#1\rrbracket}
\newcommand{\NN}{\mathds{N}}
\newcommand{\diag}{\mathop{\mathrm{diag}}}
\newcommand{\sv}{\mathbf{s}}
\newcommand{\Xv}{\mathbf{X}}
\newcommand{\EE}{\mathds{E}}
\newcommand{\pv}{\mathbf{p}}

\begin{document}
\begin{center}
\red{\textbf{This assignment is a work in progress and will be updated.}}
\end{center}
\begin{center}
\large{\textbf{CSCI-GA 3033: Honors Machine Learning} \\ Assignment 3 (Draft)\\ Due: December 7, 2026.\\ Submit on Gradescope.}
\end{center}
\begin{center}
Submit your writeup in pdf and all source code in a zip file (with proper documentation). Write a script for each programming exercise so that the TAs can easily run and verify your results. Make sure your code runs!

[Text in square brackets are hints that can be ignored.]
\end{center}
\begin{exercise}[CNN Implementation (8 pts)]
	\blue{\textbf{Note}: Please mention your Python version (and maybe the version of all other packages). For this exercise, you may submit a Python notebook.}
	
	% In this exercise you are going to run some experiments involving CNNs. You need to know \href{https://www.python.org/}{\magenta{Python}} and install the following libraries: \href{https://keras.io/}{\magenta{Keras}}, \href{https://www.tensorflow.org/install/}{\magenta{Tensorflow}}, \href{http://www.numpy.org/}{\magenta{Numpy}} and all their dependencies. You can find detailed instructions and tutorials for each of these libraries on the respective websites. 
	% [To install, try \textsf{pip install keras}. For Tensorflow, follow the installation steps on its webpage.]
	
  In this exercise you are going to run some experiments involving CNNs. You need \href{https://www.python.org/}{\magenta{Python}}, \href{https://matplotlib.org/}{\magenta{matplotlib}}, and either \href{https://pytorch.org/get-started/locally/}{\magenta{PyTorch}},  \href{https://www.tensorflow.org/tutorials/quickstart/beginner}{\magenta{TensorFlow}}, or \href{https://jax.readthedocs.io/en/latest/notebooks/quickstart.html}{\magenta{JAX}} (PyTorch is the most popular right now and is thus strongly recommended). 
  Be sure to install all necessary dependencies. You can find detailed instructions and tutorials for each of these libraries on the respective websites. 
	
	For all experiments, running on CPU is sufficient. 
  You do not need to run the code on GPUs.
  You may find \href{https://colab.research.google.com/}{Colab} useful for running things more efficiently (for free), but note that heavy usage may result in them taking away your GPU (unless you pay).
	% Before start, we suggest you review what we learned about each layer in CNN, and read at least this \href{https://keras.io/getting-started/sequential-model-guide/}{\magenta{tutorial}}.
  Before start, we suggest you review what we learned about each layer in CNN, and read the documentation and tutorial for your framework of choice (either PyTorch, TensorFlow, or JAX).
	
  \begin{enumerate}[label=\alph*)]
    \item Train a VGG11 net on the MNIST dataset (which should be loadable via, e.g., torchvision). 
		VGG11 was an earlier version of VGG16 and can be found as model A in Table 1 of this \href{https://arxiv.org/pdf/1409.1556.pdf}{\magenta{paper}}, whose Section 2.1 also gives you all the details about each layer.
		The goal is to get the loss as close to 0 loss as possible. Note that our input dimension is different from the VGG paper. You need to resize each image in MNIST from its original size $28 \times 28$ to $32 \times 32$ [make sure you understand why this is].
		
		For your convenience, we list the details of the VGG11 architecture here.
		The convolutional layers are denoted as \texttt{Conv(number of input channels, number of output channels, kernel size, stride, padding)};
		the batch normalization layers  are denoted as \texttt{BatchNorm(number of channels)};
		the max-pooling layers are denoted as \texttt{MaxPool(kernel size, stride)};
		the fully-connected layers are denoted as \texttt{FC(number of input features, number of output features)};
		the drop out layers are denoted as \texttt{Dropout(dropout ratio)}:
		\begin{verbatim}
		- Conv(001, 064, 3, 1, 1) - BatchNorm(064) - ReLU - MaxPool(2, 2)
		- Conv(064, 128, 3, 1, 1) - BatchNorm(128) - ReLU - MaxPool(2, 2)
		- Conv(128, 256, 3, 1, 1) - BatchNorm(256) - ReLU 
		- Conv(256, 256, 3, 1, 1) - BatchNorm(256) - ReLU - MaxPool(2, 2)
		- Conv(256, 512, 3, 1, 1) - BatchNorm(512) - ReLU 
		- Conv(512, 512, 3, 1, 1) - BatchNorm(512) - ReLU - MaxPool(2, 2)
		- Conv(512, 512, 3, 1, 1) - BatchNorm(512) - ReLU 
		- Conv(512, 512, 3, 1, 1) - BatchNorm(512) - ReLU - MaxPool(2, 2)
		- FC(0512, 4096) - ReLU - Dropout(0.5) 
		- FC(4096, 4096) - ReLU - Dropout(0.5) 
		- FC(4096, 10)
		\end{verbatim}
		You should use the cross-entropy loss at the end.
		
		[This experiment will take up to 1 hour on a CPU, so please be cautious of your time. If this running time is not bearable, you may cut the training set to 1/10, so only have $\sim$600 images per class instead of the regular $\sim$6000.]
		
		\item Once you've done the above, the next goal is to inspect the training process. Create the following plots:
      \begin{enumerate}[label=(\roman*)]
			\item (1 pt) test accuracy vs the number of epochs (say 3 $\sim$ 5)
			\item (1 pt) training accuracy vs the number of epochs
			\item (1 pt) test loss vs the number of epochs
			\item (1 pt) training loss vs the number of epochs
		\end{enumerate}
		[If running more than 1 epoch is computationally infeasible, simply run 1 epoch and try to record the accuracy/loss after every few minibatches.]
		
		\ans{}
		
		\item Then, it is time to inspect the generalization properties of your final model. Flip and blur the \red{test set images} using any Python library of your choice, and complete the following:		
      \begin{enumerate}[label=(\roman*)]
			\item (1 pt) test accuracy vs type of flip. Try the following two types of flipping: flip each image from left to right, and from top to bottom. Report the test accuracy after each flip. What is the effect? Please explain the effect in one sentence.
			
			As one resource for those working in PyTorch, you can read this \href{https://pytorch.org/vision/stable/transforms.html}{\magenta{doc}} to learn how to build a complex transformation pipeline. We suggest the following command for performing flipping: 
			\begin{verbatim}
			torchvision.transforms.RandomHorizontalFlip(p=1)
			torchvision.transforms.RandomVerticalFlip(p=1)
			\end{verbatim}
			
			\ans{}
			
			\item (1 pt) test accuracy vs Gaussian noise. Try adding standard Gaussian noise to each test image with variance 0.01, 0.1, 1 and report the test accuracies. What is the effect? Please explain the effect in one sentence.
			
			Again for those working in PyTorch, one way of approaching this is to apply a user-defined lambda as a new transform t which adds Gaussian noise with variance say 0.01: 
			\begin{verbatim}
			t = torchvision.transforms.Lambda(lambda x : x + 0.1*torch.randn_like(x))
			\end{verbatim}
			
			\ans{}
		\end{enumerate} 
		
		\item (2 pts) Lastly, let us verify the effect of regularization. Retrain your model with data augmentation and test again as in part 3. Report the test accuracies and explain what kind of data augmentation you use in retraining.
		
		\ans{}
		
	\end{enumerate}
	
\end{exercise}

\Needspace{8\baselineskip}
\begin{exercise}[GANs (5 pts)]
  Code provided for this exercise is assuming a PyTorch solution, you may change it as necessary if you prefer to use TensorFlow or JAX.
  You may find \href{https://pytorch.org/tutorials/beginner/dcgan_faces_tutorial.html}{this tutorial} for GANs helpful. 
  \begin{enumerate}[label=\alph*]
    \item (4 pts) Complete the implementation of a GAN in gan.py. 
      In your solution writeup, include the two produced graphs (corresponding to the training and test sets), with the epoch number on x-axis and the losses (both generator and discriminator) on the y-axis.
      Further include the produced samples from every 10th epoch (after epochs 10, 20, 30, 40, and 50).


      \ans{} 
    \item (1 pt) Describe your produced results and experiences training the GAN in this exercise.

      \ans{} 
  \end{enumerate}

\end{exercise}
\end{document}
